% ===========================================================
%  第 5 章 线性方程组 —— 高等代数【深化】拓展框（主控撰写）
%  对应 OUTLINE.md §6「深化清单」5.1 / 5.2 / 5.3
%  用法：\begin{fdeep}{标题} 灰底框；由主控插入正文相应考点旁
% ===========================================================

% ---------- 5.1 解空间 = 核空间（对应 OUTLINE 5.1）----------
\begin{fdeep}{齐次解集就是核空间：基础解系的 $n-r$ 从哪来}
齐次方程组 $Ax=0$（$A$ 为 $m\times n$）的全部解记作
\fml{N(A)=\{x\in\mathbb{R}^{n}\mid Ax=0\}}
它远不只是一个"解的集合"，而是 $\mathbb{R}^{n}$ 的一个\textbf{子空间}（叫\textbf{零空间}或\textbf{核空间}）：
因为若 $Ax_1=0$、$Ax_2=0$，则 $A(x_1+x_2)=0$、$A(kx_1)=0$，即解对加法与数乘封闭。

\medskip
\textbf{于是"基础解系"有了准确的身份}：它就是 $N(A)$ 的\textbf{一组基}。
"基础解系含 $n-r(A)$ 个向量"这句话，正是第 4 章\textbf{秩-零化度定理}的直接推论：
\fml{\dim N(A)=n-r(A)}
这里要看清三件事：
\begin{itemize}[leftmargin=2.2em,itemsep=0.25em]
\item $n$ 是\textbf{未知数个数}（不是方程个数）——很多同学错写成 $m$；
\item 之所以是 $n-r$ 而不是 $m-r$：被"压成零"的自由度取决于\textbf{输入空间}的维度 $n$；
\item 基础解系\textbf{不唯一}（任何一组基都行），但\textbf{向量个数唯一}，恒为 $n-r$。
      这也解释了为什么不同解法写出的基础解系长得不一样，却都对。
\end{itemize}
【反哺点】：服务考纲〔线性方程组〕"理解齐次线性方程组的基础解系、通解及解空间的概念"。
考场上凡是问"基础解系有几个向量""解空间维数是多少"，答案永远是 $n-r(A)$，\textbf{不必解方程就能答}。
反过来，若题目给了解空间维数，也能反推 $r(A)$。把"解空间"当成一个几何对象（子空间）而不是一堆向量，
是理解第 6 章向量空间的最佳入口。
\end{fdeep}

% ---------- 5.2 非齐次解集的结构（对应 OUTLINE 5.2）----------
\begin{fdeep}{非齐次解集：一个特解加上核空间的"平移"}
设 $Ax=b$（$b\ne0$）有解，取它的任一个解 $\eta^{*}$（称\textbf{特解}），则全部解为
\fml{\{x\mid Ax=b\}=\eta^{*}+N(A)=\{\eta^{*}+\xi\mid \xi\in N(A)\}}
\textbf{一句话理解}：$Ax=b$ 的解集是把核空间 $N(A)$ 整体"平移"了 $\eta^{*}$ 得到的。
两个解之差必是齐次解，这一点用一个等式就能看出：
\fml{A(x_1-x_2)=b-b=0\;\Longrightarrow\;x_1-x_2\in N(A)}

\textbf{为什么它"不是子空间"——这是常考的辨析点}：
$N(A)$ 含零向量，而 $b\ne0$ 时 $x=0$ 不满足 $Ax=b$，故\textbf{非齐次解集不过原点，对加法与数乘都不封闭}。
例如取两个解相加：$A(x_1+x_2)=2b\ne b$，已不再是解。所以它是\textbf{仿射子空间}（子空间平移），
只有当 $b=0$ 时才退化成子空间本身。

\medskip
\textbf{由此得到通解的"标准写法"}：
\fml{x=\eta^{*}+k_1\xi_1+k_2\xi_2+\cdots+k_{n-r}\xi_{n-r}\qquad(k_i\in\mathbb{R})}
其中 $\xi_1,\dots,\xi_{n-r}$ 是 $N(A)$ 的一组基（即基础解系），$r=r(A)$。
\textbf{要写通解必须先"找到"一个特解}——通常是令自由变量全取 $0$ 得到的最省事的那个。

【反哺点】：服务考纲〔线性方程组〕"理解非齐次线性方程组解的结构及通解的概念"。
知道"解集 = 特解 + 核空间平移"之后，$\eta^{*}+k_1\xi_1+\cdots$ 这个形状就不再是死记的公式，
而是"几何图像的文字化"。这类题的两个高频变式——"求满足 $A\eta=b$ 且与某向量正交的解"
"在解集中找一个范数最小的解"——本质都是在同一个仿射集里挑点。
\end{fdeep}

% ---------- 5.3 有解判别（对应 OUTLINE 5.3）----------
\begin{fdeep}{$A x=b$ 何时有解：$b$ 是否落在 $A$ 的列空间里}
判别法是这样的：设 $A$ 为 $m\times n$，则
\fml{Ax=b\ \text{有解}\iff r(A)=r\bigl([A\mid b]\bigr)}
\textbf{为什么要这么判——几何本质比公式好记}：
$A$ 的各列记作 $\alpha_1,\dots,\alpha_n$，方程组就是
\fml{x_1\alpha_1+x_2\alpha_2+\cdots+x_n\alpha_n=b}
也就是说：\textbf{$Ax=b$ 有解 $\iff$ $b$ 可以表示成 $A$ 的列向量的线性组合}
$\iff$ \textbf{$b$ 落在 $A$ 的列向量组张成的空间里}。
而"添上 $b$ 之后秩不变"正是"$b$ 没有带来新的维度"的形式化说法——它已经在列空间里了。

由此立刻得到两种"无解"：
\begin{itemize}[leftmargin=2.2em,itemsep=0.25em]
\item $r(A)<r([A\mid b])$：$b$ 落到了列空间外面，\textbf{无解}；
\item 若 $r(A)=r([A\mid b])$，记这个公共值为 $r$，则 $r=n$ 时\textbf{有唯一解}；$r<n$ 时\textbf{有无穷多解}（含 $n-r$ 个自由参数）。
\end{itemize}

【反哺点】：服务考纲〔线性方程组〕"理解…非齐次线性方程组有解的充分必要条件"。
把判别式理解成"$b$ 在不在列空间里"，有三个好处：
① 不会记反（谁等于谁）；② 立刻明白"无解"不是算不出来，而是几何上不可能；
③ 与第 4 章"矩阵的秩等于列向量组的秩"接上，形成一条完整的推理链，
而不是两个互不相干的结论。
\end{fdeep}
